\documentclass[12pt]{article}
\usepackage{amsmath, amssymb, graphicx, hyperref, geometry}
\geometry{a4paper, margin=1in}

\title{\textbf{RNNs from Scratch: Part 2 - Building the Architecture}}
\author{Build in Public Series}
\date{\today}

\begin{document}
\maketitle

\begin{center}
    \includegraphics[width=\textwidth]{Post-2-Banner.png}
\end{center}

\section{Introduction}
In Part 1, we explored the mathematics of recurrence. Now we implement RNNs in PyTorch, building from individual cells to complete sequence models.

\section{The RNN Cell}
Our RNNCell class implements the core recurrence:
\begin{verbatim}
h_new = tanh(W_xh @ x + W_hh @ h_prev + b)
\end{verbatim}

Key design choices:
\begin{itemize}
    \item Separate weight matrices for input and hidden projections
    \item Xavier initialization to stabilize training
    \item Optional bias terms
\end{itemize}

\section{Multi-Layer RNNs}
Stacking RNN cells creates deeper representations:
\begin{itemize}
    \item Layer 1 processes the raw input sequence
    \item Layer 2 processes the hidden states from Layer 1
    \item Each additional layer learns higher-level temporal patterns
\end{itemize}

\section{Architecture Variants}

\subsection{Sequence Classifier}
Many-to-one architecture:
\begin{verbatim}
Input:  (batch, seq_len, input_size)
Output: (batch, num_classes)
\end{verbatim}
Uses the final hidden state for classification.

\subsection{Sequence Tagger}
Many-to-many architecture:
\begin{verbatim}
Input:  (batch, seq_len, input_size)
Output: (batch, seq_len, num_tags)
\end{verbatim}
Outputs a prediction at each time step.

\subsection{CharRNN}
Character-level language model:
\begin{itemize}
    \item Embedding layer for character tokens
    \item Multi-layer RNN for sequence processing
    \item Output projection to vocabulary logits
    \item Autoregressive generation via sampling
\end{itemize}

\subsection{Bidirectional RNN}
Processes sequences in both directions:
\begin{itemize}
    \item Forward RNN: Past to future
    \item Backward RNN: Future to past
    \item Concatenate both representations
\end{itemize}

This captures context from both directions, crucial for tasks like named entity recognition.

\section{Implementation Details}

\subsection{Batch-First vs Time-First}
We use batch-first tensors $(B, T, D)$ for compatibility with PyTorch conventions, but internally process time step by time step.

\subsection{Dropout Regularization}
Applied between RNN layers (not across time) to prevent overfitting:
\begin{verbatim}
if dropout > 0:
    self.dropout = nn.Dropout(dropout)
\end{verbatim}

\subsection{Hidden State Initialization}
Hidden states are initialized to zero at the start of each sequence:
\begin{equation}
    h_0 = \mathbf{0}
\end{equation}

\section{Training Strategy}
We train on synthetic sequence tasks:
\begin{itemize}
    \item Sequence classification: Predict label from entire sequence
    \item Cross-entropy loss
    \item Adam optimizer with learning rate scheduling
\end{itemize}

\section{Conclusion}
With our RNN implemented, we can now train it and visualize how hidden states evolve. Part 3 explores hidden state dynamics and analyzes what the network learns.

\end{document}
